% TOP_PROBLEM: {"id": 30007075, "title": "A uniform pluricanonical Iitaka degree", "category_id": 5, "category": {"id": 5, "name": "algebraic_geometry", "display_name": "Algebraic Geometry", "description": "Geometric objects defined by polynomial equations.", "slug": "algebraic-geometry", "order_index": 5, "created_at": "2026-07-31T15:26:25.671Z"}, "status": "open"}
% FIRST_POSTED: 2026-10-06
% ATTEMPT_STATUS: solved
% OPENAI_CLAIM: {"schema_version":1,"source_repo":"https://github.com/openai/math","source_commit":"adc7f1241b42e322a6451854ab7e4b4c146bf78a","release_date":"2026-10-06","match":"full","resolution":"proved","independently_reviewed":false,"second_pass":"agreed","families":[{"family":"034","title":"Log abundance for compact Kähler spaces under logarithmic Iitaka subadditivity","subject":"Algebraic and complex geometry","manuscripts":[{"title":"Uniform Pluricanonical Iitaka Fibrations","pdf_path":"preprints/Uniform-Pluricanonical-Iitaka-Fibrations-October-3-2026/paper.pdf","readme_path":"preprints/Uniform-Pluricanonical-Iitaka-Fibrations-October-3-2026/README.md","date":"2026-10-03","theorem_ref":"thm:main (sections/introduction.tex:19)","statement_sources":[{"path":"preprints/Uniform-Pluricanonical-Iitaka-Fibrations-October-3-2026/build/sections/introduction.tex","line":19,"label":"thm:main"}]}],"lean":null,"reasoning_trace":null}]}
\documentclass[11pt]{article}
\usepackage[margin=25mm]{geometry}
\usepackage{fontspec}
\setmainfont{Latin Modern Roman}
\usepackage{amsmath,amssymb,mathtools}
\usepackage{unicode-math}
\setmathfont{Latin Modern Math}
\usepackage{enumitem}
\usepackage{xurl,hyperref,bookmark}
\hypersetup{colorlinks=true,urlcolor=blue}
\setcounter{secnumdepth}{0}
\setlength{\parindent}{0pt}
\setlength{\parskip}{5pt}
\setlength{\emergencystretch}{3em}
\begin{document}
\section*{A uniform pluricanonical Iitaka degree}

\subsection{OpenAI's claim}
A degree depending only on dimension is claimed to define every characteristic-zero smooth projective Iitaka fibration with nonnegative Kodaira dimension and to recover its whole function field. The cited manuscript argues for that result; its relation to this problem is specified below.

\subsection{Mathematical statement of the claimed result}
\paragraph{Mathematical claim 1.} For every $d\ge1$, OpenAI claims a positive integer $m(d)$ valid for every smooth integral projective $d$-fold $X$ over every algebraically closed characteristic-zero field with $\kappa(X)\ge0$, such that
\[
|m(d)K_X|\ne\varnothing,\qquad
k(s/t:\ 0\ne s,t\in H^0(X,m(d)K_X))=K(K_X).
\]
The field on the right is generated by all equal-degree positive pluricanonical section ratios. Thus the one complete system defines the Iitaka fibration, not merely a map of the correct image dimension.

\subsection{Outline of the argument}
Manuscript 1 (family 034, 2026-10-03): Degeneration of individual generic-fiber factors bounds their volume-form characters independently of semistable ramification. Weak positivity and flattening supply compatible small-volume bounds on the base; Frobenius-diagonal estimates and bounded curve chains then complete the dimension induction.

\subsection{Scope relative to this problem}
Manuscript 1 (family 034, 2026-10-03): The problem selects the mathematical statement summarized here, including its stated hypotheses and parameter range. The cited theorem gives a claimed proof of that assertion in full.

\subsection{Formal verification}
Family 034 has no Lean scope page in this snapshot; no formalization of these manuscript results is listed.

This site has not executed the Lean files or checked their dependencies. These descriptions report the scope of the linked declarations, rather than a new verification of the proofs.

\subsection{Status caveat}
OpenAI attributes these results to an internal model. Its release notice cautions that results lacking full formalization can contain mistakes. This site has not independently assessed the proofs; a solved label records the scope of the reported claim.

\subsection{Sources}
\begin{itemize}
\item OpenAI, \emph{Uniform Pluricanonical Iitaka Fibrations}, 2026-10-03; citation key \texttt{OAI:Uniform-Pluricanonical-Iitaka-Fibrations-October-3-2026}. \href{https://github.com/openai/math/blob/adc7f1241b42e322a6451854ab7e4b4c146bf78a/preprints/Uniform-Pluricanonical-Iitaka-Fibrations-October-3-2026/paper.pdf}{PDF (GitHub)}; \href{https://github.com/openai/math/raw/adc7f1241b42e322a6451854ab7e4b4c146bf78a/preprints/Uniform-Pluricanonical-Iitaka-Fibrations-October-3-2026/paper.pdf}{PDF (direct)}; \href{https://github.com/openai/math/blob/adc7f1241b42e322a6451854ab7e4b4c146bf78a/preprints/Uniform-Pluricanonical-Iitaka-Fibrations-October-3-2026/README.md}{Manuscript page and citation}; \href{https://github.com/openai/math/blob/main/preprints/Uniform-Pluricanonical-Iitaka-Fibrations-October-3-2026/paper.pdf}{Latest version}; Statement source: \href{https://github.com/openai/math/blob/adc7f1241b42e322a6451854ab7e4b4c146bf78a/preprints/Uniform-Pluricanonical-Iitaka-Fibrations-October-3-2026/build/sections/introduction.tex\#L19}{thm:main}.
\item \href{https://github.com/openai/math/blob/adc7f1241b42e322a6451854ab7e4b4c146bf78a/CONTENTS.md}{OpenAI catalogue entry}.
\end{itemize}
\end{document}
